\documentclass{article}
\usepackage{graphicx} % Required for inserting images
\usepackage{amsmath}
\title{Solutions to Worksheet \#1}
\author{}
\date{August 3, 2026}

\begin{document}

\maketitle
\begin{enumerate}
\item \textbf{Solution.}
We use the action in which $\pi(x_i)=x_{\pi(i)}$.  Write the three
monomials of $f$ as
\[
A=x_1x_2^2x_4^4,\qquad
B=x_2^2x_3^4x_4,\qquad
C=x_1^4x_2^2x_3.
\]
The variable $x_2$ is the only variable that has exponent $2$ in every
monomial.  Therefore, any permutation that fixes $f$ must fix $2$.

The cycle $(1\,4\,3)$ fixes $2$ and sends
\[
A\longmapsto B,\qquad B\longmapsto C,\qquad C\longmapsto A.
\]
Thus $(1\,4\,3)$ fixes the sum $f$.  Its inverse $(1\,3\,4)$ also fixes
$f$, as does the identity.

There are no other possibilities.  Indeed, after fixing $2$, the image of
$A$ must be one of $A$, $B$, or $C$.  Because the exponents $1$, $0$, and
$4$ on $x_1,x_3,x_4$ are distinct, each choice determines the permutation
completely, and the three choices give precisely the permutations above.
Hence
\[
\boxed{\{\pi\in S_4:\pi(f)=f\}
 =\{\mathrm{id},(1\,4\,3),(1\,3\,4)\}.}
\]

\item \textbf{Solution.}
Start with the monomial $x_1^7x_3^2$ and take its alternating sum over
$S_3$:
\[
f=\sum_{\pi\in S_3}\operatorname{sgn}(\pi)\,
       \pi(x_1^7x_3^2).
\]
Since the exponents $7$, $2$, and $0$ are distinct, the six monomials in
this sum are distinct.  Expanding gives
\[
\begin{aligned}
f={}&x_1^7x_3^2-x_2^7x_3^2-x_1^2x_3^7-x_1^7x_2^2\\
   &\quad+x_1^2x_2^7+x_2^2x_3^7.
\end{aligned}
\]
In particular, $x_1^7x_3^2$ occurs with coefficient $1$.

To verify the required property, let $\sigma\in S_3$.  Reindexing the sum
with $\tau=\sigma\pi$ gives
\[
\begin{aligned}
\sigma(f)
 &=\sum_{\pi\in S_3}\operatorname{sgn}(\pi)\,
      (\sigma\pi)(x_1^7x_3^2)\\
 &=\operatorname{sgn}(\sigma)
   \sum_{\tau\in S_3}\operatorname{sgn}(\tau)\,
      \tau(x_1^7x_3^2)
 =\operatorname{sgn}(\sigma)f.
\end{aligned}
\]

\medskip
\noindent\textbf{Alternative solution.}
The same polynomial can be written as the determinant
\[
f=
\begin{vmatrix}
x_1^7 & x_2^7 & x_3^7\\
1     & 1     & 1\\
x_1^2 & x_2^2 & x_3^2
\end{vmatrix}.
\]
Permuting the variables permutes the columns of the determinant, so the
determinant is multiplied by the sign of the permutation.  Expanding the
determinant also shows that the coefficient of $x_1^7x_3^2$ is $1$.

\item \textbf{Solution.}
A monomial symmetric polynomial is the sum of all distinct monomials
obtained by permuting the indicated exponents among the available
variables.  Therefore,
\[
m_{3,2}(x_1,x_2)=x_1^3x_2^2+x_1^2x_2^3,
\]
and
\[
\begin{aligned}
m_{3,2}(x_1,x_2,x_3)
={}&x_1^3x_2^2+x_1^3x_3^2+x_2^3x_1^2\\
 &+x_2^3x_3^2+x_3^3x_1^2+x_3^3x_2^2.
\end{aligned}
\]
The partition $(3,3,1,1)$ has four positive parts, so it cannot be placed
in only three variables.  Hence
\[
m_{3,3,1,1}(x_1,x_2,x_3)=0.
\]
With four variables, choose the two variables that receive exponent $3$;
the other two receive exponent $1$.  Thus
\[
\begin{aligned}
m_{3,3,1,1}(x_1,x_2,x_3,x_4)
={}&x_1^3x_2^3x_3x_4+x_1^3x_3^3x_2x_4
   +x_1^3x_4^3x_2x_3\\
 &+x_2^3x_3^3x_1x_4+x_2^3x_4^3x_1x_3
   +x_3^3x_4^3x_1x_2.
\end{aligned}
\]

\item \textbf{Solution.}
The coefficient of $x^n$ in the product is obtained by the Cauchy product:
\[
\begin{aligned}
f(x)g(x)
 &=\sum_{n=0}^{\infty}
   \left(\sum_{\substack{i+j=n\\i,j\geq 0}}j\right)x^n\\
 &=\sum_{n=0}^{\infty}
   \left(\sum_{j=0}^{n}j\right)x^n
 =\sum_{n=0}^{\infty}\frac{n(n+1)}{2}x^n.
\end{aligned}
\]
Equivalently, as formal power series,
\[
f(x)=\frac{1}{1-x},
\qquad
g(x)=\frac{x}{(1-x)^2},
\]
so
\[
\boxed{f(x)g(x)=\frac{x}{(1-x)^3}
=\sum_{n=0}^{\infty}\frac{n(n+1)}{2}x^n.}
\]

\item \textbf{Solution.}
Distributing the two sums gives
\[
\begin{aligned}
fg
 &=\left(\sum_{i=1}^{\infty}x_ix_{i+1}\right)
   \left(\sum_{j=1}^{\infty}x_j\right)\\
 &=\sum_{i=1}^{\infty}\sum_{j=1}^{\infty}
   x_ix_{i+1}x_j.
\end{aligned}
\]
Separating the cases $j=i$, $j=i+1$, and
$j\notin\{i,i+1\}$ gives the more explicit form
\[
\boxed{
fg=\sum_{i=1}^{\infty}
   \left(x_i^2x_{i+1}+x_ix_{i+1}^2\right)
 +\sum_{i=1}^{\infty}
   \sum_{\substack{j\geq 1\\j\ne i,i+1}}
   x_ix_{i+1}x_j.}
\]
The last double sum should not be simplified by discarding repetitions.
For example, $x_ix_{i+1}x_{i+2}$ occurs twice: once using the adjacent
pair $(i,i+1)$ and once using the adjacent pair $(i+1,i+2)$.

\item \textbf{Solution.}
By definition,
\[
m_{2,1}=\sum_{i\ne j}x_i^2x_j
=\sum_{i<j}\left(x_i^2x_j+x_ix_j^2\right).
\]
On the other hand,
\[
\begin{aligned}
m_2m_1
 &=\left(\sum_i x_i^2\right)\left(\sum_jx_j\right)\\
 &=\sum_i x_i^3+\sum_{i\ne j}x_i^2x_j
 =m_3+m_{2,1}.
\end{aligned}
\]
Therefore,
\[
\boxed{m_2m_1\ne m_{2,1};\qquad
m_2m_1=m_3+m_{2,1}.}
\]

\item \textbf{Solution.}
In a product $m_3m_2m_1$, choose one variable for the exponent $3$, one
for the exponent $2$, and one for the exponent $1$.  We classify the
result according to which choices are equal.
\begin{itemize}
\item If all three choices are the same, we obtain $m_6$.
\item If only the choices for exponents $3$ and $2$ are the same, we
obtain $m_{5,1}$.
\item If only the choices for exponents $3$ and $1$ are the same, we
obtain $m_{4,2}$.
\item If only the choices for exponents $2$ and $1$ are the same, we
obtain $m_{3,3}$.  Each monomial of this type occurs twice, because
either of its two variables can come from the $m_3$ factor.
\item If all three choices are distinct, we obtain $m_{3,2,1}$.
\end{itemize}
Consequently,
\[
\boxed{m_3m_2m_1
=m_6+m_{5,1}+m_{4,2}+2m_{3,3}+m_{3,2,1}.}
\]

\item \textbf{Solution.}
Recall that
\[
m_{1^r}=\sum_{\substack{A\subseteq\{1,2,\ldots\}\\|A|=r}}
\prod_{i\in A}x_i.
\]
Fix one monomial of type $(2,1^{a+b})$.  It uses $a+b+1$ distinct
variables: one variable is squared and the other $a+b$ variables occur
to the first power.

To obtain this monomial from
$m_{1^{a+1}}m_{1^{b+1}}$, the squared variable must be selected in both
factors.  Of the remaining $a+b$ variables, choose $a$ to join it in the
first factor; the remaining $b$ must be selected in the second factor.
There are
\[
\binom{a+b}{a}=\binom{a+b}{b}
\]
ways to do this.  By symmetry every monomial in
$m_{2,1^{a+b}}$ has the same coefficient.  Therefore the requested
coefficient is
\[
\boxed{\binom{a+b}{a}.}
\]

\item \textbf{Solution.}
\begin{enumerate}
\item In forming the product, we choose either $1$ or $x_i$ from each
factor.  A term obtained by choosing $r$ distinct variables is a monomial
of type $(1^r)$.  Hence
\[
\boxed{\prod_{i=1}^{\infty}(1+x_i)
=\sum_{r=0}^{\infty}m_{1^r}
=1+m_1+m_{1,1}+m_{1,1,1}+\cdots.}
\]
\item The same argument applies, but choosing $-x_i$ from $r$ factors
contributes the sign $(-1)^r$.  Therefore,
\[
\boxed{\prod_{i=1}^{\infty}(1-x_i)
=\sum_{r=0}^{\infty}(-1)^r m_{1^r}
=1-m_1+m_{1,1}-m_{1,1,1}+\cdots.}
\]
\end{enumerate}

\item \textbf{Solution.}
For each variable, use the geometric series
\[
1+x_i+x_i^2+\cdots=\frac{1}{1-x_i}.
\]
When the products are multiplied out, every monomial
$x_1^{a_1}x_2^{a_2}\cdots$ of finite total degree occurs exactly once.
Grouping monomials whose positive exponents form the same partition
$\lambda$ gives
\[
\prod_{i=1}^{\infty}\frac{1}{1-x_i}
=1+\sum_{n=1}^{\infty}\sum_{\lambda\vdash n}m_\lambda.
\]
The initial $1$ corresponds to choosing $1$ from every factor, or
equivalently to the empty partition.  Since the worksheet specifies
partitions of positive integers, the requested sum is therefore
\[
\boxed{\sum_{n=1}^{\infty}\sum_{\lambda\vdash n}m_\lambda
=\prod_{i=1}^{\infty}\frac{1}{1-x_i}-1.}
\]
If the empty partition is included and its monomial symmetric function is
defined to be $1$, then the answer is simply
$\prod_{i=1}^{\infty}(1-x_i)^{-1}$.

\end{enumerate}

\end{document}
